\chapter{Implementação}
\label{chap:implementacao}

\section{Ambiente de Simulação}
\label{sec:simulacao}

Seguindo a prática estabelecida de validar algoritmos de navegação em simulação
ROS~+~Gazebo antes de sua implantação no robô físico \cite{takaya2016}, o
ambiente de simulação utiliza o Gazebo Harmonic --- versão de suporte estendido
pareada com o ROS~2 Jazzy --- e o modelo TurtleBot4 Standard com RPLIDAR S2 (LiDAR
2D de 360°) e câmera OAK-D. O mundo \textit{warehouse} disponível no pacote
\texttt{nav2\_\allowbreak minimal\_\allowbreak tb4\_\allowbreak sim} representa um armazém industrial com prateleiras,
barreiras e espaços abertos, oferecendo diversidade de geometria para testar o SLAM
em diferentes condições de observabilidade.

A escolha do \texttt{nav2\_\allowbreak minimal\_\allowbreak tb4\_\allowbreak sim} em detrimento do
\texttt{turtlebot4\_\allowbreak gz\_\allowbreak bringup} resultou de problemas de compatibilidade entre os
plugins \texttt{irobot\_\allowbreak create\_\allowbreak gz\_\allowbreak plugins} --- projetados para Ignition Fortress ---
e o Gazebo Harmonic, que renomeou a API de mensagens de \texttt{ignition.msgs} para
\texttt{gz.msgs}. A versão minimal contorna esse problema utilizando o plugin de
DiffDrive nativo do Gazebo Harmonic com bridges corretos, incluindo o topic
\texttt{cmd\_\allowbreak vel} como \texttt{geometry\_\allowbreak msgs/\allowbreak msg/\allowbreak Twist} (não \texttt{TwistStamped}).
Essa distinção exigiu que os parâmetros de navegação do Nav2 fossem gerados em tempo
de launch com \texttt{enable\_\allowbreak stamped\_\allowbreak cmd\_\allowbreak vel: false} em todos os onze nós afetados,
via leitura e patch programático do arquivo \texttt{nav2.yaml} padrão.

\section{O Orquestrador em Detalhe}
\label{sec:impl_orquestrador}

O orquestrador (\texttt{fleet\_\allowbreak orchestrator}) é implementado como um nó ROS~2 em Python
com \texttt{MultiThreadedExecutor}. A estrutura de estado por robô, \texttt{RobotRuntime},
mantém seis campos: \texttt{nav\_\allowbreak state}, \texttt{is\_\allowbreak recording}, \texttt{is\_\allowbreak navigating},
\texttt{route\_\allowbreak poses}, \texttt{last\_\allowbreak saved} e \texttt{pending\_\allowbreak goal\_\allowbreak pose}. Este último
é o mais crítico para a gravação: ao receber o callback de resultado do action Nav2,
se o status for \texttt{STATUS\_\allowbreak SUCCEEDED} e o robô estiver em modo de gravação,
a pose enviada como meta é imediatamente appendada à lista de rotas --- pois a pose
real confirmada pelo Nav2 é a melhor estimativa disponível da posição atingida.

A lógica de gravação baseada em polling de TF (\texttt{\_\allowbreak record\_\allowbreak tick}) funciona em
paralelo para ambientes onde o Nav2 não reporta SUCCEEDED (e.g., quando o controlador
para dentro da tolerância antes de publicar o callback). O tick a 5\,Hz consulta a
cadeia TF e adiciona a pose ao buffer se o robô deslocou ao menos 10\,cm ou girou
mais de 5° desde a última amostra, evitando redundância em trechos estacionários.


\subsection{Exemplo: os Serviços \texttt{start\_\allowbreak record} e \texttt{stop\_\allowbreak record}}
\label{sec:impl_servicos_exemplo}

Para tornar concreta a discussão anterior, o Código~\ref{lst:start_stop_record}
reproduz, tal como publicado no repositório público do projeto, os \textit{callbacks}
que implementam o início e o fim de uma gravação de rota. Note-se a rejeição explícita
de requisições inconsistentes com o estado corrente do robô (\texttt{ALREADY\_\allowbreak NAVIGATING})
e a limpeza do buffer de poses e do último ponto salvo (\texttt{last\_\allowbreak saved = None}) a
cada novo início de gravação --- sem essa limpeza, uma segunda gravação herdaria pontos
da anterior, corrompendo silenciosamente o arquivo de rota resultante.

\begin{lstlisting}[style=pycode, caption={Callbacks \texttt{\_\allowbreak cb\_\allowbreak start\_\allowbreak record} e \texttt{\_\allowbreak cb\_\allowbreak stop\_\allowbreak record} (\texttt{fleet\_\allowbreak orchestrator/\allowbreak main.py}).}, label={lst:start_stop_record}]
def _cb_start_record(self, req, resp):
    rid = req.robot_id
    if not self._known_robot(rid):
        resp.success = False
        resp.message = f"Unknown robot {rid!r}"
        resp.error_code = "UNKNOWN_ROBOT"
        return resp
    st = self._state[rid]
    if st.is_navigating:
        resp.success = False
        resp.message = "Cannot start recording while navigating"
        resp.error_code = "ALREADY_NAVIGATING"
        return resp
    st.route_poses.clear()
    st.last_saved = None
    st.is_recording = True
    st.nav_state = "recording"
    st.current_route = req.route_name
    st.last_error = ""
    resp.success = True
    resp.message = "Recording started"
    resp.error_code = ""
    return resp

def _cb_stop_record(self, req, resp):
    rid = req.robot_id
    if not self._known_robot(rid):
        resp.success = False
        resp.message = f"Unknown robot {rid!r}"
        resp.error_code = "UNKNOWN_ROBOT"
        return resp
    st = self._state[rid]
    route_name = st.current_route
    st.is_recording = False
    n = len(st.route_poses)
    if n == 0:
        st.nav_state = "idle"
        st.current_route = ""
        resp.success = True
        resp.message = "Route length: 0"
        resp.error_code = ""
        return resp
    path = self._save_route_yaml(rid, route_name, st.route_poses)
    st.route_poses.clear()
    st.last_saved = None
    st.nav_state = "idle"
    st.current_route = ""
    resp.success = True
    resp.message = f"Saved {n} poses to {path}"
    resp.error_code = ""
    return resp
\end{lstlisting}

Um segundo aspecto relevante do orquestrador é a separação entre os papéis
experimentais e a permissão de movimento, aplicada em \texttt{\_\allowbreak cb\_\allowbreak play\_\allowbreak route}
e \texttt{\_\allowbreak cb\_\allowbreak go\_\allowbreak to\_\allowbreak point} por meio do método \texttt{\_\allowbreak motion\_\allowbreak allowed},
que consulta o mapa de papéis carregado de \texttt{roles.yaml} e retorna
\texttt{False} para qualquer robô que não seja \texttt{MUUT}. Requisições de
movimento para um FUUT recebem o código de erro \texttt{ROLE\_\allowbreak NOT\_\allowbreak ALLOWED}
antes mesmo de qualquer tentativa de contato com o action server do Nav2,
evitando gastar o \textit{timeout} de \texttt{wait\_\allowbreak for\_\allowbreak server} em uma
requisição que já se sabe inválida por construção.

\subsection{A Ponte REST--ROS~2 no Backend}
\label{sec:impl_backend_ponte}

O backend não utiliza um cliente ROS~2 nativo para encaminhar comandos de
serviço vindos da interface web: em vez disso, cada endpoint HTTP dispara
um processo \texttt{ros2 service call} via \texttt{subprocess}, como mostra
o Código~\ref{lst:backend_bridge}. A justificativa para essa escolha, aparentemente
pouco ortodoxa, está detalhada na Seção~\ref{sec:backend_web}: evitar que o
processo do backend precise ele mesmo inicializar um nó \texttt{rclpy} completo
só para chamadas pontuais de serviço, mantendo essa inicialização (via
\texttt{rclpy.init()} em uma \textit{thread} dedicada) restrita à publicação
contínua de status, mapa e pose, que de fato exige um nó ROS~2 residente.

\begin{lstlisting}[style=pycode, caption={Função \texttt{\_\allowbreak run\_\allowbreak ros2\_\allowbreak service}, usada por todos os endpoints \texttt{/\allowbreak api/\allowbreak *} que encaminham comandos ao orquestrador (\texttt{backend/\allowbreak main.py}).}, label={lst:backend_bridge}]
def _run_ros2_service(srv, srv_type, request_json, timeout=10):
    cmd = (
        f"source /opt/ros/jazzy/setup.bash 2>/dev/null; "
        f"source {ROS_WS}/install/setup.bash 2>/dev/null; "
        f"ros2 service call {srv} {srv_type} '{request_json}'"
    )
    try:
        r = subprocess.run(
            ["bash", "-c", cmd],
            env=_ros_env(),
            capture_output=True,
            text=True,
            timeout=timeout,
            cwd=ROS_WS,
        )
        out = (r.stdout or "").strip() + (r.stderr or "").strip()
        if r.returncode != 0:
            return False, out or "ros2 service call failed"
        return True, out
    except subprocess.TimeoutExpired:
        return False, "timeout"
    except Exception as e:
        return False, str(e)
\end{lstlisting}

Essa abordagem tem um custo mensurável: cada chamada paga o overhead de
inicializar um processo \texttt{ros2 service call} (tipicamente
50--150\,ms na máquina de desenvolvimento), o que a torna inadequada para
comandos de alta frequência, mas aceitável para o padrão de uso da interface
--- cliques esparsos do pesquisador, não um laço de controle. O preço em
latência é trocado por uma superfície de acoplamento mínima entre o backend
e o \textit{workspace} ROS~2: o backend só precisa saber o nome e o tipo do
serviço, nunca importar diretamente os módulos gerados de \texttt{fleet\_\allowbreak msgs}
para essas chamadas específicas.

\section{Pipeline de Coleta e Análise}
\label{sec:impl_pipeline}

O script \texttt{experiment\_\allowbreak repeatability.py} implementa o protocolo descrito no
Capítulo~\ref{chap:metodologia} como uma sequência de chamadas a serviços ROS~2.
Sua estrutura em dois subcomandos (\texttt{record} e \texttt{replay}) permite compor
qualquer protocolo experimental a partir de blocos elementares, seja pela linha de
comando ou pela interface web.

A análise offline em \texttt{analyze\_\allowbreak runs.py} lê os bags via \texttt{rosbag2\_\allowbreak py},
desserializando diretamente as mensagens de tipos conhecidos
(\texttt{nav\_\allowbreak msgs/\allowbreak Odometry}, \texttt{geometry\_\allowbreak msgs/\allowbreak PoseWithCovarianceStamped})
a partir do formato CDR armazenado no MCAP. A versão do script utilizada na
campanha final desta dissertação calcula o RMSE pairwise após reamostragem em
tempo normalizado: cada trajetória é projetada para uma grade comum no intervalo
$[0,1]$, e as coordenadas $x$ e $y$ são interpoladas nessa grade antes do cálculo
do erro quadrático médio. Essa escolha evita associar pontos apenas pelo índice
da amostra, o que poderia introduzir artefatos quando duas execuções percorrem
a mesma rota com pequenas diferenças de velocidade.

Para fins de auditoria, o próprio \texttt{summary.json} produzido pela análise
registra o modo de reamostragem usado. Na campanha final
\texttt{dissertation\_\allowbreak clean01\_\allowbreak final\_\allowbreak manual}, o campo
\texttt{resampling.mode} indica \texttt{time}, com 100 amostras na grade
normalizada.


\subsection{Reamostragem e RMSE: o Código que Sustenta a Métrica}
\label{sec:impl_rmse_codigo}

O Código~\ref{lst:rmse_impl} reproduz as funções centrais de reamostragem e
cálculo de RMSE em \texttt{analyze\_\allowbreak runs.py}. A função
\texttt{\_\allowbreak normalised\_\allowbreak time\_\allowbreak grid} converte o vetor temporal de cada bag
para o intervalo $[0,1]$; em seguida, \texttt{numpy.interp} calcula posições
homólogas na mesma grade normalizada. A implementação por índice permanece
disponível como opção de compatibilidade (\texttt{--resample-mode index}), mas
não foi usada na campanha final reportada no Capítulo~\ref{chap:resultados}.

\begin{lstlisting}[style=pycode, caption={Reamostragem temporal normalizada e RMSE pairwise (\texttt{scripts/\allowbreak analyze\_\allowbreak runs.py}).}, label={lst:rmse_impl}]
def _normalised_time_grid(t):
    if len(t) < 2:
        return np.zeros(len(t), dtype=np.float64)
    duration = float(t[-1] - t[0])
    if duration <= 1e-9:
        return np.linspace(0.0, 1.0, len(t), dtype=np.float64)
    return (t - t[0]) / duration

def _resample_xy_by_time(t, xy, samples):
    if len(xy) < 2 or samples < 2:
        return xy
    tn = _normalised_time_grid(t)
    unique_t, unique_idx = np.unique(tn, return_index=True)
    unique_xy = xy[unique_idx]
    grid = np.linspace(0.0, 1.0, samples, dtype=np.float64)
    return np.column_stack([
        np.interp(grid, unique_t, unique_xy[:, 0]),
        np.interp(grid, unique_t, unique_xy[:, 1]),
    ])

def _rmse(a, b):
    return float(np.sqrt(np.mean(np.sum((a - b) ** 2, axis=1))))
\end{lstlisting}

A consequência prática dessa escolha é que duas trajetórias percorridas em
velocidades ligeiramente distintas continuam comparáveis, pois os pontos
homólogos representam a mesma fração normalizada do percurso temporal de cada
execução. Essa formulação não elimina todas as ameaças à validade --- diferenças
muito grandes de duração ainda podem indicar que as execuções não são
experimentalmente equivalentes ---, mas reduz o artefato geométrico introduzido
pela reamostragem puramente por índice.

\section{Interface Web}
\label{sec:impl_web}

A interface web é uma aplicação React com backend FastAPI. O mapa SLAM é renderizado
em um canvas HTML5 com coordenadas alinhadas ao sistema de referência ROS: a origem
do canvas corresponde ao canto inferior esquerdo do mapa, e o eixo Y é invertido
(positivo para cima em ROS, negativo no canvas). A conversão utiliza os metadados
do \texttt{OccupancyGrid} \cite{elfes1989} (\texttt{resolution}, \texttt{origin\_\allowbreak x}, \texttt{origin\_\allowbreak y})
publicados pelo SLAM Toolbox.

O mapa pré-construído do \textit{warehouse} (\texttt{warehouse.pgm} +
\texttt{warehouse.yaml}) é utilizado como fundo estático com alinhamento perfeito:
como ambos o mapa SLAM ao vivo e o PGM utilizam o mesmo sistema de coordenadas
do mundo Gazebo, a sobreposição é pixel-a-pixel. O pesquisador vê paredes, prateleiras
e o robô em tempo real, podendo clicar para definir waypoints diretamente na planta.

\section{Considerações Finais}
\label{sec:impl_conclusao}

As decisões de implementação descritas neste capítulo --- o patch runtime do
\texttt{nav2.yaml}, o \texttt{MultiThreadedExecutor} para evitar deadlocks, a
desserialização CDR sem tipo registrado, a conversão de coordenadas canvas--mundo
--- resultaram de problemas concretos encontrados durante o desenvolvimento.
Cada uma representa uma contribuição de engenharia que outros grupos de pesquisa
que adotarem o framework não precisarão replicar. O próximo capítulo apresenta
o arranjo de avaliação que valida o sistema resultante.
